Curating a training dataset from a mass of unfiltered data is a necessary and
influential step in any large-scale machine learning pipeline. Deciding how to
curate such a dataset is a challenging problem that has attracted substantial
recent
interest~\citep{fang2022data,abbas2023semdedup,engstrom2024dsdm,gadre2024datacomp}.
In this section, we frame pre-training data selection as an optimization
problem, and then solve this problem by first-order optimizing with
metagradients. Applying our method to the DataComp-small benchmark
\citep{gadre2024datacomp}, we greatly improve on the state-of-the-art (our
  improvement over state-of-the-art is roughly the same as the
improvement of state-of-the-art over training on random data).

\subsubsection{Setup}
The goal of {dataset selection} is to choose a training data subset (out of
a broad pool of data) that maximizes trained machine learning model
performance.
Given this goal,
dataset selection has a natural interpretation as a
combinatorial metaparameter optimization problem.
In particular, in the language of Section \ref{subsec:meta_gradient},
for a training set of size $n$, let
\begin{itemize}
  \item[(a)] the metaparameters $\mathbf{c} \in \mathcal{C} :=
    \mathbb{Z}_{\geq 0}^n$ be non-negative data counts representing
    the number of times each training sample repeats in the training data;
  \item[(b)] the algorithm $\mathcal{A}$ be a standard large-scale
    learning procedure, which runs on a training set
    comprising $c_i$ copies of each sample $i$ for $i \in [n]$;
  \item[(c)] the output function $\phi$ be the loss of the trained model on a target
    distribution $\dtarg{}$.
\end{itemize}
Then, defining $f(\mathbf{c}) := \phi(\mathcal{A}(\mathbf{c}))$ (as
in Section~\ref{subsec:meta_gradient}),
our goal is to find the data counts $\mathbf{c^*}$ that solve
\begin{equation}
  \label{eq:ds}
  \mathbf{c}^* := \argmin_{\mathbf{c} \in \mathcal{C}} f(\mathbf{c}).
\end{equation}

\subsubsection{Gradient descent on training data}
\label{sec:gdod}
Metagradients let us \textit{directly} minimize the target task
loss \eqref{eq:ds}
with respect to the choice of training data.
At a high level, our algorithm operates as follows:
we start with a randomly chosen set of training data, then
iteratively update the dataset selection using
metagradients with respect to importance weights
placed on each training datapoint.
The specifics of our method are in Algorithm~\ref{alg:depsing};
we describe its core ideas below.

\paragraph{Idea 1: A surrogate algorithm.} We cannot use metagradients to
optimize \eqref{eq:ds} directly, because the metaparameters of
interest $\mathbf{c}$ are
discrete counts (and so the algorithm $\mathcal{A}$ is non-differentiable with 
respect to $\mathbf{c}$). 
To circumvent this problem, we relax $\mathcal{A}$: we define a surrogate algorithm
$\mathcal{A}_\mathbf{c}'$ that takes in a {\em continuous}
metaparameter $\mathbf{z} \in
\mathbb{R}^n$, whose metagradient we \textit{can} compute, then optimize using
the metagradient on $\mathcal{A}_\mathbf{c}'$.

This surrogate learning algorithm $\mathcal{A}_\mathbf{c}'$ maps a metaparameter 
$\mathbf{z} \in \mathbb{R}^n$ (representing a perturbation to training data weights) 
to a machine learning model.
The surrogate is defined by a set
of counts $\mathbf{c} \in \mathbb{Z}_+^n$, and a hyperparameter $k$ denoting a
specific training iteration, both of which we bake into the surrogate algorithm
itself. 
Given a metaparameter $\mathbf{z} \in \mathbb{R}^n$, the algorithm
$\mathcal{A}_\mathbf{c}'$ trains a model ``as usual'' using the fixed counts
$\mathbf{c}$. That is, it makes $c_i$ copies of each training sample $i$,
shuffles and partitions the data into batches, and then at each iteration
minimizes the batch loss with a step---just as the original learning algorithm
$\mathcal{A}$. At iteration $k$, however, in addition to the original loss on
the $k$-th batch, the algorithm upweights {\em each} training sample $i$
according to the metaparameter $z_i$. In other words, the objective at iteration
$t$ of the surrogate algorithm $\mathcal{A}_\mathbf{c}'$ is
\begin{equation*}
  \label{eq:surrogate_objective}
  \ell'_t(\theta) :=
  \begin{cases}
    \sum_{x \in t^{\text{th}} \text{ batch}} \ell(x; \theta) &
    \!\!\text{if } t \neq k \\
    \sum_{x \in t^{\text{th}} \text{ batch}} \ell(x; \theta) +
    \sum_{i=1}^n z_i \ell(x_i; \theta) &\!\!\text{if } t = k
  \end{cases}
\end{equation*}
where $\ell(x; \theta)$ is the training loss on example $x$.

Observe that when $\mathbf{z} = \mathbf{0}_n$, the algorithm
$\mathcal{A}_\mathbf{c}'$  is identical to the standard learning algorithm
$\mathcal{A}$. And while $\mathcal{A}$ was a function of (nondifferentiable)
discrete data counts $\mathbf{c}$, $\mathcal{A}_\mathbf{c}'$ is differentiable
with respect
its input $\mathbf{z}$, and so we can compute the metagradient
\[
  \mathbf{g} := \nabla_{\mathbf{z}}
  \phi\big(\mathcal{A}_\mathbf{c}'(\mathbf{z})\big)\big\vert_{\mathbf{z}
  = \mathbf{0}_n}.
\]
Intuitively, the entries of the metagradient $\mathbf{g}$ capture the effect of
adding an infinitesimal amount of each training sample $i$ to the training data
at iteration $k$. A positive entry $g_i$ indicates that adding an infinitesimal
amount of sample $i$ to the training data would increase the loss,  and a
negative entry indicates that adding an infinitesimal amount of sample $i$ to
the training data would decrease the loss; the slot at $i$ represents the
(estimated) effect of adding a copy of sample $i$ to the training
data at every batch containing the sample.

\paragraph{Idea 2: Block coordinate descent.}
We then use the metagradient $\mathbf{g}$ to iteratively update our selected
dataset. We update data counts as
\begin{equation}
  \label{eq:update}
  \mathbf{c} \gets \mathbf{c} - \mathrm{sign}(\mathbf{g}) \odot
  \mathbf{m}, \quad \mathbf{m} \sim \mathrm{Bernoulli}(p),
\end{equation}
where $p$ is a hyperparameter controlling the fraction of sample counts to
update. This algorithm resembles a block coordinate descent algorithm
\citep{ortega2000iterative}, with the main difference being that we take signed
gradient steps with step size $1$ (projected onto non-negative integers) to
ensure that the counts remain well-defined. As a result, $p$ implicitly controls
the algorithm's step size.

Applying \eqref{eq:update} concludes a single optimization step.
By repeating this process of estimating the metagradient, updating our
counts vector, then constructing a new training dataset, we iteratively
improve the selected data.  Pseudocode for our algorithm can be found in Algorithm \ref{alg:depsing}.

\begin{algorithm}[h]
    \SetKwFor{For}{\textcolor{forcolor}{for}}{\textcolor{forcolor}{do}}{}
    \SetKw{Return}{\textcolor{forcolor}{Return}}{}
    \SetKwFunction{traversereverse}{\textcolor{fncolor}{reverse\_inorder\_traversal}}
    \DontPrintSemicolon
	\KwIn{initial data counts $\mathbf{c} \in \mathbb{Z}_{\geq 0}^n$, learning algorithm $\mathcal{A}$, output function $\phi$}
    \nonl\hspace{0pt}{\textbf{Hyperparameters:}} step size $p$, \# opt steps $T$, iteration number $k$ \;
    \For{$t \gets 1$ \color{forcolor}{\textbf{to}} \color{black}$T$}{
		$\mathbf{z} \gets \mathbf{0}_n$ \tcp{Build input to surrogate}
		$\mathbf{g} \gets \frac{\partial \phi(\mathcal{A}'_\mathbf{c}(\mathbf{z}))}{\partial \mathbf{z}}$ \tcp{Calculate metagradient using \textsc{Replay}}
		$\mathbf{m} \gets$ sample from $\textrm{Bernoulli}(p)$ \tcp{Sample indices to step on}
		$\mathbf{c} \gets \mathbf{c} - \mathrm{sign}(\mathbf{g}) \odot \mathbf{m}$ \tcp{Take optimization step}
    }
	\Return $\mathbf{c}$ \tcp{Return final data counts}
	\caption{Dataset selection using using metagradient descent (MGD).}
    \label{alg:depsing}
\end{algorithm}

\subsubsection{Results}
We evaluate our data selection algorithm using DataComp
\cite{gadre2024datacomp}, a standardized framework for evaluating data selection
methods for multimodal models. Algorithm~\ref{alg:depsing} greatly
improves on the state-of-the-art for the benchmark. Below, we
describe the setting, outline our
method, and conclude with our results.

\paragraph{Setting.}
DataComp \citep{gadre2024datacomp} is a multimodal model training competition 
and benchmark for evaluating dataset selection methods. 
DataComp provides a \textit{fixed} learning algorithm
chosen in advance by the organizers and a large fixed \emph{candidate pool} of
internet data. The goal is to choose a subset of the
candidate pool---possibly with repeated datapoints---that yields the best-performing
model after training with the given learning algorithm, as measured by a
predetermined set of 38 benchmarks. Given a submission subset, the mean score on
the evaluation datasets for a model trained with that subset is taken as the
final ``score.''  DataComp offers four separate ``scales'' requiring different
amounts of compute; we focus on the \emph{small} scale in this paper
due to compute limitations.

\paragraph{Method.}
We select data with MGD (Algorithm~\ref{alg:depsing}) to minimize loss on
data on a ``target set'' that is distributionally similar to the DataComp
benchmark tasks, and select hyperparameters with a held-out ``validation
set.'' In particular, we construct target and validation sets by taking
samples from the DataComp evaluation tasks with extra samples available beyond
those used in the DataComp test set (e.g., ImageNet, one of the tasks in
DataComp, has a training set in addition to the test set evaluated in DataComp).
See Appendix~\ref{app:clip_details} for the exact details of the target and
validation sets, the precise hyperparameters used with
Algorithm~\ref{alg:depsing}, and a discussion on scalability (including further
engineering details on executing our algorithm efficiently).

\paragraph{Results.}
MGD greatly outperforms the current state-of-the-art: the difference in accuracy
between MGD and the current best method is roughly as large as the difference
between the previous state-of-the-art (EcoDatum~\citep{ecodatum2024ecodatum})
and training on randomly chosen data
(cf.\ Figure~\ref{fig:datacomp_score_over_time}).
Inspecting scores over the course of the optimization in Figure
\ref{fig:datacomp_score_over_time}, we find that only a few steps are necessary
to outperform previous methods.

\begin{figure}[ht]
  \centering
  \begin{subfigure}[t]{0.49\textwidth}
    \raisebox{-\height}{\pgfplotsset{colormap/Set2}
\pgfplotsset{scaled y ticks=false}
\pgfplotsset{grid style={dashed,gray}}

\begin{tikzpicture}
    \begin{axis}[
      xlabel={Metagradient steps},
      ylabel={DataComp Score},
      height={4.8cm},
      width={7.5cm},
      colormap name={Set2},
      ytick distance=0.03,
      xmin=-3,
      xmax=57,
      grid=both
    ]

    \addplot[mark=none, index of colormap=3, domain=-10:60, samples=2, dashed, very thick] {0.1377730386384371};
    \addlegendentry{No filtering}
    \addplot[mark=none, index of colormap=1, domain=-10:60, samples=2, dashed, very thick] {0.173};
    \addlegendentry{Best DataComp baseline}
    \addplot[mark=none,, index of colormap=2, domain=-10:60, samples=2, dashed, very thick] {0.182};
    \addlegendentry{Prev. SOTA (EcoDatum)}
    \addplot [scatter,
                index of colormap=0,
                very thick,
                scatter/@pre marker code/.code={3
                    \def\markopts{index of colormap=0}
                    \expandafter\scope\expandafter[\markopts]
                },
                scatter/@post marker code/.code={
                    \endscope
                },] 
    table[
      x=it,
      y={score},
      col sep=comma,
    ]{pgffigs/datacomp_small_result/datacomp_small_result.csv};
    \addlegendentry{MGD-DS (ours)}

    \legend{}
    \end{axis}
\end{tikzpicture}
}
  \end{subfigure}
  \hfill
  \begin{subfigure}[t]{0.5\textwidth}
    \begin{tabular}[t]{@{}lcc@{}}
\toprule
\textbf{Method}  &  \textbf{Score} & $\Delta$ \\
\midrule
\tikz[baseline=-0.5ex] {\draw [dashed, color=pgfcolor3, very thick] (0,0) -- (0.6,0)} Baseline: No filtering  & 0.13 & -- \\
\tikz[baseline=-0.5ex] {\draw [dashed, color=pgfcolor1, very thick] (0,0) -- (0.6,0)} Best baseline from \cite{gadre2024datacomp}  & 0.17 & \textcolor{mygreen}{+0.04} \\
\tikz[baseline=-0.5ex] {\draw [dashed, color=pgfcolor2, very thick] (0,0) -- (0.6,0)} Previous SOTA \cite{ecodatum2024ecodatum}  & 0.18 & \textcolor{mygreen}{+0.05} \\
\tikz[baseline=-0.5ex] {\draw [color=pgfcolor0, very thick] (0,0) -- (0.6,0); \filldraw [color=pgfcolor0] (0.3,0) circle (0.08);} \textbf{MGD-DS (ours)} & \textbf{0.22} & \textbf{\textcolor{mygreen}{+0.09}} \\
\bottomrule
\end{tabular}

  \end{subfigure}
  \caption{MGD dataset selection greatly outperforms existing methods (improving
      over the previous SOTA by as much as the previous SOTA improves over no
    filtering at all). We compare DataComp scores for MGD (over
    optimization steps),
    training on the entire candidate pool, the best baseline
    originally proposed by
  DataComp, and the previous SOTA~\citep{ecodatum2024ecodatum}.}
  \label{fig:datacomp_score_over_time}
\end{figure}

